\chapter{2010年新课标卷（理）}

\section{单选题}

\begin{problem}
已知集合 \(A = \{x \mid |x| \leqslant 2, x \in \R\}\)，\(B = \{x \mid \sqrt{x} \leqslant 4, x \in \Z\}\)，则 \(A \cap B =\)（\quad）

\choices
    {\((0,2)\)}
    {\([0,2]\)}
    {\(\{0,2\}\)}
    {\(\{0,1,2\}\)}
\end{problem}

\begin{answer}
D.
\end{answer}

\begin{problem}
已知复数 \(z = \frac{\sqrt{3} + \i}{(1 - \sqrt{3}\i)^2}\)，\(\overline{z}\) 是 \(z\) 的共轭复数，则 \(z \cdot \overline{z} =\)（\quad）

\choices
    {\(\frac{1}{4}\)}
    {\(\frac{1}{2}\)}
    {\(1\)}
    {\(2\)}
\end{problem}

\begin{answer}
A.
\end{answer}

\begin{problem}
曲线 \( y = \frac{x}{x + 2} \) 在点 \((-1, -1)\) 处的切线方程为（\quad）

\choices
    {\( y = 2x + 1 \)}
    {\( y = 2x - 1 \)}
    {\( y = -2x - 3 \)}
    {\(y = -2x - 2\)}
\end{problem}

\begin{answer}
A.
\end{answer}

\begin{problem}
如图，质点 \(P\) 在半径为 \(2\) 的圆周上逆时针运动，其初始位置为 \(P_{0}(\sqrt{2}, - \sqrt{2})\)，角速度为 \(1\)，那么点 \(P\) 到 \(x\) 轴的距离 \(d\) 关于时间 \(t\) 的函数图象大致为（\quad）

\bitmapfigure[max width=0.5\linewidth]{img/2010/new_curriculum_science/q4_fig1.png}

\choices
    {\choicebitmap[width=0.15\paperwidth]{img/2010/new_curriculum_science/q4_optA.png}}
    {\choicebitmap[width=0.15\paperwidth]{img/2010/new_curriculum_science/q4_optB.png}}
    {\choicebitmap[width=0.15\paperwidth]{img/2010/new_curriculum_science/q4_optC.png}}
    {\choicebitmap[width=0.15\paperwidth]{img/2010/new_curriculum_science/q4_optD.png}}
\end{problem}

\begin{answer}
C.
\end{answer}

\begin{problem}
已知命题 \(p_{1}\)：函数 \( y = 2^{x} - 2^{-x} \) 在 \( \R \) 上为增函数；\(p_{2}\)：函数 \( y = 2^{x} + 2^{-x} \) 在 \( \R \) 上为减函数. 则在命题 \(q_{1}: p_{1} \vee p_{2}\)，\(q_{2}: p_{1} \wedge p_{2}\)，\(q_{3}: (\neg p_{1}) \vee p_{2}\) 和 \(q_{4}: p_{1} \wedge (\neg p_{2})\) 中，真命题是（\quad）

\choices
    {\(q_{1}\)，\(q_{3}\)}
    {\(q_{2}\)，\(q_{3}\)}
    {\(q_{1}\)，\(q_{4}\)}
    {\(q_{2}\)，\(q_{4}\)}
\end{problem}

\begin{answer}
C.
\end{answer}

\begin{problem}
某种种子每粒发芽的概率都为 \(0.9\)，现播种了 \(1000\text{ 粒}\). 对于没有发芽的种子，每粒需再补种 \(2\text{ 粒}\)，补种的种子数记为 \(X\)，则 \(X\) 的数学期望为（\quad）

\choices
    {\(100\)}
    {\(200\)}
    {\(300\)}
    {\(400\)}
\end{problem}

\begin{answer}
B.
\end{answer}

\begin{problem}
如果执行如图所示的框图，输入 \(N = 5\)，则输出的数等于（\quad）

\bitmapfigure[float=inline,max width=7.5cm]{img/2010/new_curriculum_science/q7_fig1.png}

\choices
    {\(\frac{5}{4}\)}
    {\(\frac{4}{5}\)}
    {\(\frac{6}{5}\)}
    {\(\frac{5}{6}\)}
\end{problem}

\begin{answer}
B.
\end{answer}

\begin{problem}
设偶函数 \( f(x) \) 满足 \( f(x) = x^3 - 8  \)（\(x \geqslant 0\)），则 \( \{x \mid f(x - 2) > 0\} =\)（\quad）

\choices
    {\(\{x\mid x < -2 \text{或} x > 4\}\)}
    {\(\{x\mid x < 0 \text{或} x > 4\}\)}
    {\(\{x\mid x < 0 \text{或} x > 6\}\)}
    {\(\{x\mid x < -2 \text{或} x > 2\}\)}
\end{problem}

\begin{answer}
B.
\end{answer}

\begin{problem}
若 \(\cos \alpha = -\frac{4}{5}\)，\(\alpha\) 是第三象限的角，则 \(\frac{1 + \tan \frac{\alpha}{2}}{1 - \tan \frac{\alpha}{2}} =\)（\quad）

\choices
    {\(-\frac{1}{2}\)}
    {\( \frac{1}{2} \)}
    {\(2\)}
    {\(-2\)}
\end{problem}

\begin{answer}
A.
\end{answer}

\begin{problem}
设三棱柱的侧棱垂直于底面，所有棱长都为 \(a\)，顶点都在一个球面上，则该球的表面积为（\quad）

\choices
    {\(\pi a^2\)}
    {\(\frac{7}{3}\pi a^2\)}
    {\(\frac{11}{3}\pi a^2\)}
    {\(5\pi a^{2}\)}
\end{problem}

\begin{answer}
B.
\end{answer}

\begin{problem}
已知函数 \( f(x) = \begin{cases}
|\lg x|, & 0 < x \leqslant 10,\\
- \frac{1}{2} x + 6, & x > 10,
\end{cases} \) 若 \(a\)，\(b\)，\(c\) 互不相等，且 \( f(a) = f(b) = f(c) \)，则 \( abc \) 的取值范围是（\quad）

\choices
    {\((1,10)\)}
    {\((5,6)\)}
    {\((10,12)\)}
    {\((20,24)\)}
\end{problem}

\begin{answer}
C.
\end{answer}

\begin{problem}
已知双曲线 \(E\) 的中心为原点，\(F(3,0)\) 是 \(E\) 的焦点，过 \(F\) 的直线 \(l\) 与 \(E\) 相交于 \(A\)，\(B\) 两点，且 \(AB\) 的中点为 \(N(-12, -15)\)，则 \(E\) 的方程为（\quad）

\choices
    {\(\frac{x^2}{3} - \frac{y^2}{6} = 1\)}
    {\(\frac{x^2}{4} - \frac{y^2}{5} = 1\)}
    {\(\frac{x^2}{6} - \frac{y^2}{3} = 1\)}
    {\(\frac{x^2}{5} - \frac{y^2}{4} = 1\)}
\end{problem}

\begin{answer}
B.
\end{answer}

\section{填空题}

\begin{problem}
设 \( y = f(x) \) 为区间 \([0,1]\) 上的连续函数，且恒有 \( 0 \leqslant f(x) \leqslant 1 \)，可以用随机模拟方法近似计算积分 \(\int_{0}^{1} f(x) \mathrm{d}x\). 先产生两组（每组 \(N\text{ 个}\)）区间 \([0,1]\) 上的均匀随机数 \(x_1\)，\(x_2\)，\(\cdots\)，\(x_N\) 和 \(y_1\)，\(y_2\)，\(\cdots\)，\(y_N\)，由此得到 \(N\text{ 个}\)点 \((x_i, y_i)\)（\(i = 1\)，\(2\)，\(\cdots\)，\(N\)），再数出其中满足 \( y_i \leqslant f(x_i) \)（\(i = 1\)，\(2\)，\(\cdots\)，\(N\)）的点数 \(N_1\)，那么由随机模拟方法可得积分 \( \int_0^1 f(x) \mathrm{d}x \) 的近似值为\fillinblank{}.
\end{problem}

\begin{answer}
\(\frac{N_{1}}{N}\).
\end{answer}

\begin{problem}
正视图为一个三角形的几何体可以是\fillinblank{}.（写出三种）
\end{problem}

\begin{answer}
三棱锥、三棱柱、圆锥.
\end{answer}

\begin{problem}
过点 \(A(4,1)\) 的圆 \(C\) 与直线 \(x - y - 1 = 0\) 相切于点 \(B(2,1)\)，则圆 \(C\) 的方程为\fillinblank{}.
\end{problem}

\begin{answer}
\((x-3)^2+y^2=2\).
\end{answer}

\begin{problem}
在 \(\triangle ABC\) 中，\(D\) 为边 \(BC\) 上一点，\(BD = \frac{1}{2} DC\)，\(\angle ADB = 120^{\circ}\)，\(AD = 2\)，若 \(\triangle ADC\) 的面积为 \(3 - \sqrt{3}\)，则 \(\angle BAC =\fillinblank{}\).
\end{problem}

\begin{answer}
\(60^\circ\).
\end{answer}

\section{解答题}

\begin{problem}
设数列 \(\{a_{n}\}\) 满足 \(a_{1} = 2\)，\(a_{n+1} - a_{n} = 3 \cdot 2^{2n-1}\).

\begin{enumerate}
\item 求数列 \( \{a_{n}\} \) 的通项公式；
\item 令 \( b_{n} = na_{n} \)，求数列 \( \{b_{n}\} \) 的前 \( n \) 项和 \( S_{n} \).
\end{enumerate}
\end{problem}

\begin{solution}
\begin{enumerate}
\item
由已知，
\[
a_n=a_1+\sum_{k=1}^{n-1}(a_{k+1}-a_k)
=2+\sum_{k=1}^{n-1}3\cdot 2^{2k-1}.
\]
因为 \(2^{2k-1}=\dfrac{4^k}{2}\)，所以
\[
a_n=2+\frac{3}{2}\sum_{k=1}^{n-1}4^k
=2+\frac{3}{2}\cdot\frac{4(4^{n-1}-1)}{4-1}
=2\cdot 4^{n-1}=2^{2n-1}.
\]
因此数列的通项公式为 \(a_n=2^{2n-1}\).

\item
又 \(b_n=na_n=n2^{2n-1}=\dfrac{n4^n}{2}\)，于是
\[
S_n=\frac{1}{2}\sum_{k=1}^{n}k4^k.
\]
令 \(T_n=\sum_{k=1}^{n}k4^k\)，则
\[
4T_n=\sum_{k=1}^{n}k4^{k+1}=\sum_{k=2}^{n+1}(k-1)4^k.
\]
两式相减得
\[
3T_n=n4^{n+1}-\sum_{k=1}^{n}4^k
=n4^{n+1}-\frac{4(4^n-1)}{3},
\]
从而
\[
T_n=\frac{4+(3n-1)4^{n+1}}{9},
\qquad
S_n=\frac{2+2(3n-1)4^n}{9}.
\]
所以 \(\{b_n\}\) 的前 \(n\) 项和为
\[
S_n=\frac{2\bigl((3n-1)4^n+1\bigr)}{9}.
\]
\end{enumerate}
\end{solution}

\begin{problem}
如图，已知四棱锥 \(P - ABCD\) 的底面为等腰梯形，\(AB \parallel CD\)，\(AC \perp BD\)，垂足为 \(H\)，\(PH\) 是四棱锥的高，\(E\) 为 \(AD\) 中点.

\begin{enumerate}
\item 证明：\(PE \perp BC\)；
\item 若 \(\angle APB = \angle ADB = 60^{\circ}\)，求直线 \(PA\) 与平面 \(PEH\) 所成角的正弦值.
\end{enumerate}

\bitmapfigure[max width=0.34\linewidth]{img/2010/new_curriculum_science/q18_fig1.png}

\end{problem}

\begin{solution}
\begin{enumerate}
\item
在底面所在平面内建立直角坐标系，设
\[
A(-u,0,0),\quad B(u,0,0),\quad C(v,h,0),\quad D(-v,h,0),
\]
其中 \(u>0\)，\(v>0\)，且 \(ABCD\) 为等腰梯形. 由
\[
\overrightarrow{AC}=(u+v,h,0),\qquad
\overrightarrow{BD}=(-(u+v),h,0)
\]
及 \(AC\perp BD\)，得 \(h^2-(u+v)^2=0\). 因 \(h>0\)，故 \(h=u+v\). 两条对角线的交点为
\[
H=(0,u,0)
\]
设 \(P=(0,u,p)\)，其中 \(p>0\)，则
\[
E=\left(-\frac{u+v}{2},\frac{u+v}{2},0\right)
\]
于是
\[
\overrightarrow{PE}=\left(-\frac{u+v}{2},\frac{v-u}{2},-p\right),
\qquad
\overrightarrow{BC}=(v-u,u+v,0),
\]
从而
\[
\overrightarrow{PE}\cdot\overrightarrow{BC}
=-\frac{u+v}{2}(v-u)+\frac{v-u}{2}(u+v)=0
\]
所以 \(PE\perp BC\).

\item
又 \(PH\perp\) 底面，故 \(PH\perp BC\). 由于 \(PE\) 和 \(PH\) 是平面 \(PEH\) 内的两条相交直线，所以 \(BC\perp\) 平面 \(PEH\).

由 \(\angle ADB=60^{\circ}\)，有
\[
\overrightarrow{DA}=(v-u,-u-v,0),\qquad
\overrightarrow{DB}=(u+v,-u-v,0),
\]
故
\[
\cos\angle ADB
=\frac{\overrightarrow{DA}\cdot\overrightarrow{DB}}{|DA||DB|}
=\frac{2v(u+v)}{\sqrt{2(u^2+v^2)}\,\sqrt{2}(u+v)}
=\frac{v}{\sqrt{u^2+v^2}}=\frac12
\]
所以 \(u^2=3v^2\). 又由 \(\angle APB=60^{\circ}\)，
\[
\overrightarrow{PA}=(-u,-u,-p),\qquad
\overrightarrow{PB}=(u,-u,-p),
\]
从而
\[
\frac{\overrightarrow{PA}\cdot\overrightarrow{PB}}{|PA||PB|}
=\frac{p^2}{2u^2+p^2}=\frac12
\]
即 \(p^2=2u^2\). 设直线 \(PA\) 与平面 \(PEH\) 所成的角为 \(\theta\)，由于 \(BC\) 是该平面的法向量，
\[
\sin\theta
=\frac{|\overrightarrow{PA}\cdot\overrightarrow{BC}|}{|PA||BC|}
=\frac{2uv}{\sqrt{2u^2+p^2}\,\sqrt{2(u^2+v^2)}}
\]
代入 \(p^2=2u^2\) 和 \(u^2=3v^2\)，得
\[
\sin\theta=\frac{2uv}{2u\cdot2\sqrt{2}v}
=\frac{\sqrt{2}}{4}
\]
\end{enumerate}
\end{solution}

\begin{problem}
为调查某地区老人是否需要志愿者提供帮助，用简单随机抽样方法从该地区调查了 \(500\text{ 位}\)老年人，结果如下：

\begin{center}
\small
\setlength{\tabcolsep}{0.8em}
\renewcommand{\arraystretch}{1.2}
\begin{tabular}{|c|c|c|}
\hline
\makebox[3.6cm][l]{%
  \rule[-18.82pt]{0pt}{43.12pt}%
  \smash{\rlap{\raisebox{1.05cm}{\hspace{-0.3cm}\rotatebox{-19.75}{\rule{4.49cm}{0.4pt}}}}}%
  \rlap{\raisebox{0.74cm}{\hspace{2.6cm}性别}}%
  \rlap{\raisebox{-0.4cm}{\hspace{0.12cm}是否需要志愿者}}%
}
& 男 & 女 \\
\hline
需要 & \(40\) & \(30\) \\
\hline
不需要 & \(160\) & \(270\) \\
\hline
\end{tabular}
\end{center}

\begin{enumerate}
\item 估计该地区老年人中，需要志愿者提供帮助的老年人的比例；
\item 能否有 \(99\%\) 的把握认为该地区的老年人是否需要志愿者提供帮助与性别有关？
\item 根据（2）的结论，能否提出更好的调查方法来估计该地区老年人中，需要志愿者帮助的老年人的比例？说明理由.
\end{enumerate}

附：
\begin{center}
\begin{tabular}{c|c c c}
\(P(K^2\geqslant k)\) & \(0.050\) & \(0.010\) & \(0.001\) \\
\hline
\(k\) & \(3.841\) & \(6.635\) & \(10.828\)
\end{tabular}
\end{center}

\begin{center}
\(\displaystyle K ^{2} = \frac {n (a d - b c) ^{2}}{(a + b) (c + d) (a + c) (b + d)}\)
\end{center}
\end{problem}

\begin{solution}
\begin{enumerate}
\item 样本中需要志愿者的老年人有 \(40+30=70\) 人，因此该地区老年人中需要志愿者提供帮助的比例估计为
\[
\hat p=\frac{70}{500}=0.14
\]

\item 取 \(a=40\)，\(b=30\)，\(c=160\)，\(d=270\)，\(n=500\)，则
\[
K^2=\frac{500(40\times270-30\times160)^2}
{(40+30)(160+270)(40+160)(30+270)}
=\frac{500\times6000^2}{70\times430\times200\times300}
\approx9.967
\]
在显著性水平 \(0.01\) 下，临界值为 \(6.635\). 由于 \(9.967>6.635\)，所以拒绝“是否需要志愿者与性别相互独立”的假设，有 \(99\%\) 的把握认为该地区老年人是否需要志愿者提供帮助与性别有关.

\item 可以采用按性别分层的随机抽样方法. 先将老年人按性别分为男、女两层，再在两层中分别进行随机抽样，并按各层老年人在总体中的比例加权估计总体比例. 若男性、女性总体人数分别为 \(N_1,N_2\)，样本中的相应比例估计分别为 \(\hat p_1,\hat p_2\)，则总体比例估计为
\[
\hat p=\frac{N_1\hat p_1+N_2\hat p_2}{N_1+N_2}
\]
若按总体比例分配样本量，则能使样本中男女比例与总体接近；由于性别与是否需要志愿者有关，这种分层抽样通常比简单随机抽样具有更小的抽样误差，因而是更好的调查方法.
\end{enumerate}
\end{solution}

\begin{problem}
设 \(F_{1}\)，\(F_{2}\) 分别是椭圆 \(E: \frac{x^{2}}{a^{2}} + \frac{y^{2}}{b^{2}} = 1 \) （\(a > b > 0\)）的左，右焦点，过 \(F_{1}\) 且斜率为 1 的直线 \(l\) 与 \(E\) 相交于 \(A\)，\(B\) 两点，且 \(|AF_{2}|\)，\(|AB|\)，\(|BF_{2}|\) 成等差数列.

\begin{enumerate}
\item 求 \(E\) 的离心率；
\item 设点 \(P(0, -1)\) 满足 \(|PA| = |PB|\)，求 \(E\) 的方程.
\end{enumerate}
\end{problem}

\begin{solution}
设椭圆的半焦距为 \(c\)，则 \(F_1(-c,0)\)，\(F_2(c,0)\)，且
\[
c^2=a^2-b^2
\]
直线 \(l\) 的方程为 \(y=x+c\). 设它与椭圆的两个交点的横坐标为 \(x_1,x_2\)，则将 \(y=x+c\) 代入椭圆方程，得
\[
\frac{x^2}{a^2}+\frac{(x+c)^2}{b^2}=1
\]
整理得
\[
(a^2+b^2)x^2+2a^2cx+a^2(c^2-b^2)=0
\]
其判别式为
\[
\begin{aligned}
\Delta
&=(2a^2c)^2-4(a^2+b^2)a^2(c^2-b^2)\\
&=8a^2b^4,
\end{aligned}
\]
其中用了 \(c^2=a^2-b^2\). 因此，二次方程两根之差的绝对值为
\[
|x_1-x_2|=\frac{2\sqrt{2}ab^2}{a^2+b^2}
\]
由于直线斜率为 \(1\)，所以
\[
|AB|=\sqrt{(x_1-x_2)^2+(y_1-y_2)^2}
=\sqrt{2}|x_1-x_2|
=\frac{4ab^2}{a^2+b^2}
\]

焦点 \(F_1\) 在椭圆内部且在弦 \(AB\) 上，因此 \(|AB|=|AF_1|+|BF_1|\). 由椭圆的定义，
\[
|AF_1|+|AF_2|=2a,\qquad |BF_1|+|BF_2|=2a
\]
于是
\[
|AF_2|+|BF_2|=4a-|AB|
\]
因为 \(|AF_2|,|AB|,|BF_2|\) 成等差数列，
\[
2|AB|=|AF_2|+|BF_2|=4a-|AB|
\]
故 \(|AB|=\frac{4a}{3}\). 结合弦长公式，得
\[
\frac{4ab^2}{a^2+b^2}=\frac{4a}{3},\qquad a^2=2b^2
\]
所以
\[
c^2=a^2-b^2=b^2,\qquad e=\frac ca=\frac{\sqrt{2}}{2}
\]

再设弦中点为 \(N\). 由二次方程根和，
\[
x_1+x_2=-\frac{2a^2c}{a^2+b^2},\qquad
x_N=\frac{x_1+x_2}{2}=-\frac{a^2c}{a^2+b^2}=-\frac{2c}{3}
\]
条件 \(|PA|=|PB|\) 表明 \(P(0,-1)\) 在弦 \(AB\) 的垂直平分线上. 由于 \(AB\) 的斜率为 \(1\)，该垂直平分线为 \(y=-x-1\)；它与 \(l:y=x+c\) 的交点就是 \(N\)，故
\[
x_N=-\frac{c+1}{2}
\]
于是
\[
-\frac{2c}{3}=-\frac{c+1}{2},\qquad c=3
\]
从而 \(a^2=2c^2=18\)，\(b^2=c^2=9\)，椭圆方程为
\[
\frac{x^2}{18}+\frac{y^2}{9}=1
\]
\end{solution}

\begin{problem}
设函数 \( f(x) = \e^{x} - 1 - x - ax^2 \).

\begin{enumerate}
\item 若 \(a = 0\)，求 \(f(x)\) 的单调区间；
\item 若当 \(x \geqslant 0\) 时 \(f(x) \geqslant 0\)，求 \(a\) 的取值范围.
\end{enumerate}
\end{problem}

\begin{solution}
\begin{enumerate}
\item 当 \(a=0\) 时，
\[
f(x)=\e^x-1-x,\qquad f'(x)=\e^x-1
\]
当 \(x<0\) 时，\(\e^x<1\)，故 \(f'(x)<0\)；当 \(x>0\) 时，\(\e^x>1\)，故 \(f'(x)>0\). 所以 \(f(x)\) 在 \((-\infty,0)\) 上单调递减，在 \((0,+\infty)\) 上单调递增.

\item 对任意 \(x\geqslant0\)，先证明 \(\e^x\geqslant1+x+\dfrac{x^2}{2}\).
令 \(g(x)=\e^x-1-x-\dfrac{x^2}{2}\)，再令 \(h(x)=\e^x-1-x\)，则
\[
h'(x)=\e^x-1\geqslant0,\qquad h(0)=0,
\]
所以 \(h(x)\geqslant0\)；从而
\[
g'(x)=\e^x-1-x=h(x)\geqslant0,\qquad g(0)=0,
\]
故 \(g(x)\geqslant0\)，即 \(\e^x\geqslant1+x+\dfrac{x^2}{2}\). 于是
\[
f(x)=\e^x-1-x-ax^2\geqslant\left(\frac12-a\right)x^2
\]
因此当 \(a\leqslant\dfrac12\) 时，恒有 \(f(x)\geqslant0\).

反之，若对任意 \(x\geqslant0\) 都有 \(f(x)\geqslant0\)，则对任意 \(x>0\)，
\[
a\leqslant\frac{\e^x-1-x}{x^2}
\]
令 \(x\to0^+\)，并连续两次使用洛必达法则，有
\[
\lim_{x\to0^+}\frac{\e^x-1-x}{x^2}
=\lim_{x\to0^+}\frac{\e^x-1}{2x}
=\lim_{x\to0^+}\frac{\e^x}{2}
=\frac12
\]
所以 \(a\leqslant\dfrac12\).

所以 \(a\) 的取值范围为 \(\left(-\infty,\dfrac12\right]\).
\end{enumerate}
\end{solution}

\begin{problem}
如图，已知圆上的弧 \(\widehat{AC} = \widehat{BD}\)，过 \(C\) 点的圆的切线与 \(BA\) 的延长线交于 \(E\) 点，证明：

\begin{enumerate}
\item \(\angle ACE = \angle BCD\)；
\item \(BC^{2} = BE\cdot CD\).
\end{enumerate}

\bitmapfigure[max width=0.65\linewidth]{img/2010/new_curriculum_science/q22_fig1.png}

\end{problem}

\begin{solution}
\begin{enumerate}
\item
由弧 \(\widehat{AC}=\widehat{BD}\)，所对的圆周角相等，因此
\[
\angle ABC=\angle BCD
\]
又由切线与弦所夹的角等于弦所对的圆周角，得
\[
\angle ACE=\angle ABC
\]
所以 \(\angle ACE=\angle BCD\).

\item
由于 \(B,A,E\) 共线，
\[
\angle CBE=\angle CBA=\angle BCD
\]
其中最后一个等式由弧相等得到. 又由切线与弦所夹的角定理，
\[
\angle BCE=\angle CDB
\]
所以 \(\triangle BCE\sim\triangle CDB\)，从而
\[
\frac{BC}{CD}=\frac{BE}{BC}
\]
即
\[
BC^2=BE\cdot CD
\]
\end{enumerate}
\end{solution}

\begin{problem}
已知直线 \(C_1: \begin{cases}
x = 1 + t\cos \alpha,\\
y = t\sin \alpha,
\end{cases}\) （\(t\) 为参数），圆 \(C_2: \begin{cases}
x = \cos \theta,\\
y = \sin \theta.
\end{cases}\) （\(\theta\) 为参数）.

\begin{enumerate}
\item 当 \(\alpha = \frac{\pi}{3}\) 时，求 \(C_1\) 与 \(C_2\) 的交点坐标；
\item 过坐标原点 \(O\) 作 \(C_{1}\) 的垂线，垂足为 \(A\)，\(P\) 为 \(OA\) 的中点，当 \(\alpha\) 变化时，求点 \(P\) 轨迹的参数方程，并指出它是什么曲线.
\end{enumerate}
\end{problem}

\begin{solution}
\begin{enumerate}
\item 当 \(\alpha=\dfrac{\pi}{3}\) 时，直线 \(C_1\) 上的点为
\[
\left(1+\frac{t}{2},\frac{\sqrt{3}}{2}t\right)
\]
代入圆 \(C_2:x^2+y^2=1\)，得
\[
\left(1+\frac{t}{2}\right)^2+\left(\frac{\sqrt{3}}{2}t\right)^2=1
\]
即 \(t(t+1)=0\). 当 \(t=0\) 时交点为 \((1,0)\)；当 \(t=-1\) 时交点为 \(\left(\dfrac12,-\dfrac{\sqrt{3}}2\right)\).

\item 直线 \(C_1\) 上点的坐标可写为
\[
(x,y)=(1,0)+t(\cos\alpha,\sin\alpha)
\]
垂足 \(A\) 满足 \(\overrightarrow{OA}\perp(\cos\alpha,\sin\alpha)\)，所以
\[
(1+t\cos\alpha)\cos\alpha+t\sin^2\alpha=0
\]
即 \(t=-\cos\alpha\). 因此
\[
A=(\sin^2\alpha,-\sin\alpha\cos\alpha)
\]
而 \(P\) 是 \(OA\) 的中点，故点 \(P\) 的轨迹参数方程为
\[
\begin{cases}
x=\dfrac12\sin^2\alpha,\\
y=-\dfrac12\sin\alpha\cos\alpha
\end{cases}
\]
利用二倍角公式，
\[
x-\frac14=-\frac14\cos2\alpha,\qquad y=-\frac14\sin2\alpha
\]
所以
\[
\left(x-\frac14\right)^2+y^2=\frac1{16}
\]
当 \(\alpha\) 变化时，点 \(P\) 的轨迹是圆心为 \(\left(\dfrac14,0\right)\)、半径为 \(\dfrac14\) 的圆.
\end{enumerate}
\end{solution}

\begin{problem}
设函数 \( f(x) = |2x - 4| + 1 \).

\begin{enumerate}
\item 画出函数 \( y = f(x) \) 的图象；
\item 若不等式 \( f(x) \leqslant ax \) 的解集非空，求 \( a \) 的取值范围.
\end{enumerate}

\bitmapfigure[max width=0.35\linewidth]{img/2010/new_curriculum_science/q24_fig1.png}
\end{problem}

\begin{solution}
\begin{enumerate}
\item
\[
f(x)=|2x-4|+1=
\begin{cases}
-2x+5,&x\leqslant2,\\
2x-3,&x\geqslant2
\end{cases}
\]
所以函数图象是以 \((2,1)\) 为顶点、左支斜率为 \(-2\)、右支斜率为 \(2\) 的折线，即开口向上的 \(V\) 形图象.

\item 若不等式有解，则分情况讨论.

当 \(x\leqslant2\) 时，
\[
5-2x\leqslant ax\quad\Longleftrightarrow\quad 5\leqslant(a+2)x
\]
若 \(x<0\)，要有解必须 \(a+2<0\)，即 \(a<-2\)；若 \(0<x\leqslant2\)，则由右端的最大值不小于 \(5\)，得 \(2(a+2)\geqslant5\)，即 \(a\geqslant\dfrac12\). 而 \(x=0\) 时不等式不成立.

当 \(x\geqslant2\) 时，
\[
2x-3\leqslant ax\quad\Longleftrightarrow\quad (a-2)x\geqslant-3
\]
若 \(a\geqslant2\)，显然有解；若 \(a<2\)，由于 \(x\geqslant2\)，要有解必须
\[
2(a-2)\geqslant-3
\]
即 \(a\geqslant\dfrac12\).

综上，不等式的解集非空时
\[
a\in\left(-\infty,-2\right)\cup\left[\frac12,+\infty\right)
\]
反之，若 \(a<-2\)，取 \(x=\dfrac{5}{a+2}<0\) 即可取等号；若 \(a\geqslant\dfrac12\)，取 \(x=2\) 即有 \(1\leqslant2a\). 因此该范围确实可取.
\end{enumerate}
\end{solution}
